\documentclass{article}
\usepackage{amsmath,amssymb,geometry}
\geometry{a4paper,margin=1.5cm}
\usepackage{ctex}
\usepackage{enumitem}
\usepackage{xcolor}

\newcommand{\comment}[1]{\textcolor{blue}{\small\textbf{【指导】#1}}}

\begin{document}
	
	\subsubsection*{(1) 直线 $AC$ 到平面 $PEF$ 的距离}
	
	\begin{enumerate}[label=\textbf{\arabic*.},leftmargin=*]
		\item \textbf{建系与坐标化}：以 $A$ 为原点，$AB,AD,AP$ 所在直线分别为 $x,y,z$ 轴建立空间直角坐标系。
		\[
		A(0,0,0),\; B(1,0,0),\; C(1,1,0),\; D(0,1,0),\; P(0,0,1),\;
		E\left(\frac{1}{2},0,0\right),\; F\left(1,\frac{1}{2},0\right)
		\]
		\comment{建系原则：选取两两垂直且交于一点的棱为坐标轴，这里 $AB,AD,AP$ 两两垂直。}
		
		\item \textbf{求平面 $PEF$ 的法向量 $\boldsymbol{n}$}：
		\[
		\overrightarrow{PE} = \left(\frac{1}{2},0,-1\right),\quad \overrightarrow{PF} = \left(1,\frac{1}{2},-1\right)
		\]
		设 $\boldsymbol{n}=(x,y,z)$，由 $\boldsymbol{n} \perp \overrightarrow{PE}$ 且 $\boldsymbol{n} \perp \overrightarrow{PF}$：
		\[
		\begin{cases}
			\dfrac{1}{2}x - z = 0 \\[6pt]
			x + \dfrac{1}{2}y - z = 0
		\end{cases}
		\]
		由第一式得 $z = \dfrac{1}{2}x$，代入第二式：
		\[
		x + \frac{1}{2}y - \frac{1}{2}x = 0 \Rightarrow \frac{1}{2}x + \frac{1}{2}y = 0 \Rightarrow y = -x
		\]
		取 $x=2$，得 $y=-2,\;z=1$，故 $\boldsymbol{n}=(2,-2,1)$，$|\boldsymbol{n}| = \sqrt{4+4+1}=3$。
		
		\comment{法向量坐标取整数便于后续计算，可任意取值，只需满足方程组。}
		
		\item \textbf{验证线面平行}：
		$\overrightarrow{AC} = (1,1,0)$，计算 $\overrightarrow{AC} \cdot \boldsymbol{n} = 1\times2 + 1\times(-2) + 0\times1 = 0$。
		故 $AC \parallel$ 平面 $PEF$，直线到平面的距离可转化为点到平面的距离。
		\comment{线面距离存在的前提是直线与平面平行，数量积为 $0$ 是平行的代数条件。}
		
		\item \textbf{求点 $A$ 到平面 $PEF$ 的距离}：
		取平面上一点 $E$，$\overrightarrow{AE} = \left(\dfrac{1}{2},0,0\right)$，则
		\[
		d = \frac{|\overrightarrow{AE} \cdot \boldsymbol{n}|}{|\boldsymbol{n}|}
		= \frac{\left|\dfrac{1}{2}\times2 + 0\times(-2) + 0\times1\right|}{3}
		= \frac{1}{3}
		\]
		\comment{点面距公式：$d = \dfrac{|\text{斜向量} \cdot \text{法向量}|}{|\text{法向量}|}$，斜向量需从平面上一点指向目标点。}
		
		因此，直线 $AC$ 到平面 $PEF$ 的距离为 $\dfrac{1}{3}$。
	\end{enumerate}
	
	\subsubsection*{(2) $BD$ 与平面 $PEF$ 的夹角}
	
	\begin{enumerate}[label=\textbf{\arabic*.},leftmargin=*]
		\item \textbf{直线的方向向量}：$\overrightarrow{BD} = D - B = (0-1,1-0,0-0) = (-1,1,0)$。
		
		\item \textbf{求线面角的正弦值}：
		设 $BD$ 与平面 $PEF$ 所成角为 $\theta$，则
		\[
		\sin\theta = |\cos\langle \overrightarrow{BD}, \boldsymbol{n} \rangle|
		= \frac{|\overrightarrow{BD} \cdot \boldsymbol{n}|}{|\overrightarrow{BD}| \cdot |\boldsymbol{n}|}
		= \frac{|(-1)\times2 + 1\times(-2) + 0\times1|}{\sqrt{(-1)^2+1^2+0^2} \times 3}
		= \frac{|-2-2|}{\sqrt{2} \times 3}
		= \frac{4}{3\sqrt{2}} = \frac{2\sqrt{2}}{3}
		\]
		\comment{线面角公式：$\sin\theta = |\cos\langle \text{方向向量}, \text{法向量} \rangle|$，结果取锐角。}
		
		\item \textbf{结论}：$BD$ 与平面 $PEF$ 的夹角 $\theta = \arcsin\dfrac{2\sqrt{2}}{3}$。
	\end{enumerate}
	
\end{document}